\subsection{معرفی و انگیزه}
این مکانیزم بر مبنای مبانی نظری پیشرفته در ادبیات توجه و فشرده‌سازی \cite{eckart1936approximation} توسعه یافته است.

مکانیزم توجه با فشرده‌سازی دنباله بر پایه تجزیه مقادیر منفرد \lr{(Low-Rank SVD-Compressed Sequence Attention)} با هدف غلبه بر پیچیدگی محاسباتی و حافظه‌ای درجه دو $\mathcal{O}(N^2)$ در طول دنباله‌های بلند طراحی شده است. در لایه‌های توجه کلاسیک، تشکیل ماتریس توجه تمام‌به‌تمام مستلزم ضرب داخلی میان تمامی جفت‌توکن‌های دنباله است؛ این در حالی است که فضای زیرفضای ویژگی‌های کلید و مقدار در اغلب لایه‌های ترنسفورمر دارای بعد مؤثر یا رتبه ذاتی بسیار پایینی است.

ماژول \lr{SVDLinearAttention} (ثبت‌شده با شناسه‌های \lr{svd\_attention}، \lr{svd\_seq}، \lr{svd\_sequence} و \lr{svd\_sequence\_attention}) با استفاده از تجزیه مقادیر منفرد ماتریس تصویربرداری کلید ($\mathbf{W}_K$) در زمان مقداردهی، یک زیرفضای پایدار متعامد با رتبه $r \ll N$ استخراج می‌کند. سپس با نگاشت دنباله کلید و مقدار بر روی این زیرفضای فشرده، بعد دنباله از $N$ به $r$ انقباض یافته و پیچیدگی توجه به $\mathcal{O}(N \cdot r \cdot d)$ تبدیل می‌گردد که مقیاس‌پذیری خطی نسبت به طول توالی را به ارمغان می‌آورد.
